\chapter{Puntspiegeling en parallellogrammen}\label{ch:g7:central}

Draai een figuur een halve draai om een punt: dat is de
\emph{\omterm{def:g7:central:def}{puntspiegeling}}, de tweede transformatie van dit boek na de spiegelingen van
\cref{ch:g6:symmetry}. Haar sterfiguur is het \omterm{def:g7:central:parallelogram}{parallellogram} --- de \omterm{def:g4:geometry:polygon}{vierhoek} die
symmetrisch is om het snijpunt van haar eigen diagonalen.

\section{Puntspiegeling}

\begin{definition}[Spiegelbeeld om een punt]\label{def:g7:central:def}
Het \emph{spiegelbeeld}\index{puntspiegeling} van een punt $M$ om een punt $O$ is
het punt $M'$ waarvoor $O$ het \emph{midden} van $[MM']$ is. De gespiegelde figuur
is wat het origineel wordt na een halve draai ($180^\circ$) om $O$.
\end{definition}

\begin{omfigure}
\begin{tikzpicture}[scale=1.0]
  \fill (0,0) circle (1.6pt) node[below, font=\small] {$O$};
  \fill[omDef] (-2.3,1.0) circle (1.8pt) node[above, font=\small] {$M$};
  \fill[omThm] (2.3,-1.0) circle (1.8pt) node[below, font=\small] {$M'$};
  \draw[dashed] (-2.3,1.0) -- (2.3,-1.0);
  \draw (-1.2,0.45) -- (-1.1,0.6);
  \draw (1.1,-0.6) -- (1.2,-0.45);
  \fill[omDef] (-1.4,-0.8) circle (1.8pt) node[below, font=\small] {$N$};
  \fill[omThm] (1.4,0.8) circle (1.8pt) node[above, font=\small] {$N'$};
  \draw[dashed] (-1.4,-0.8) -- (1.4,0.8);
\end{tikzpicture}

{\small Spiegelen om $O$: elk punt en zijn beeld liggen met $O$ op één \omterm{def:g6:lines:objects}{lijn}, aan
weerszijden op gelijke afstanden (streepjes).}
\end{omfigure}

\begin{method}[Het spiegelbeeld van een punt construeren]\label{met:g7:central:construct}
Om het \omterm{def:g7:central:def}{spiegelbeeld} $M'$ van $M$ om $O$ te construeren:
\begin{enumerate}
  \item teken de \omterm{def:g6:lines:objects}{lijn} $(MO)$ en verleng haar voorbij $O$;
  \item meet $MO$ (met de lat of de passer);
  \item zet $M'$ op het verlengde met $OM' = OM$.
\end{enumerate}
Op ruitjespapier: tel de horizontale en verticale stappen van $M$ naar $O$, en
herhaal \emph{dezelfde} stappen vanaf $O$ om bij $M'$ te komen.
\end{method}

\begin{example}[Op een raster]\label{ex:g7:central:grid}
Ligt $M$ $3$ vakjes links en $1$ vakje boven $O$, dan ligt zijn \omterm{def:g7:central:def}{spiegelbeeld} $M'$
$3$ vakjes \emph{rechts} en $1$ vakje \emph{onder} $O$: een \omterm{def:g7:central:def}{puntspiegeling} keert
beide richtingen tegelijk om (een spiegeling in een \omterm{def:g6:lines:objects}{lijn} keert er maar één om).
\end{example}

\begin{proposition}[Wat een halve draai bewaart]\label{prop:g7:central:preserved}
Een \omterm{def:g7:central:def}{puntspiegeling} bewaart lengtes, hoeken, omtrekken en \omterm{def:g6:measure:area}{oppervlaktes}; het beeld van
een \omterm{def:g6:lines:objects}{lijn} is een \emph{parallelle} \omterm{def:g6:lines:objects}{lijn}, het beeld van een \omterm{def:g6:lines:objects}{lijnstuk} is een \omterm{def:g6:lines:perp}{parallel}
\omterm{def:g6:lines:objects}{lijnstuk} van dezelfde lengte, en het beeld van een cirkel is een cirkel met dezelfde
\omterm{def:g3:shapes:shapes}{radius} (met het \omterm{def:g7:central:def}{spiegelbeeld} van het middelpunt als middelpunt).
\end{proposition}
\admitted

\begin{remark}
Vergelijk met de spiegelingen (\cref{prop:g6:symmetry:preserved}): beide bewaren
lengtes en hoeken; maar de spiegeling klapt figuren om (een ``b'' wordt een ``d''),
terwijl de halve draai ze leesbaar houdt --- op hun kop (een ``b'' wordt een ``q'').
En anders dan een spiegeling stuurt een halve draai elke \omterm{def:g6:lines:objects}{lijn} naar een \omterm{def:g6:lines:objects}{lijn} die
\omterm{def:g6:lines:perp}{parallel} aan haar is.
\end{remark}

\begin{definition}[Symmetriemiddelpunt]\label{def:g7:central:center}
Een punt $O$ is een \emph{symmetriemiddelpunt}\index{symmetriemiddelpunt} van een
figuur als de halve draai om $O$ de figuur precies op zichzelf afbeeldt. Zo hebben
de letters S, N en Z een symmetriemiddelpunt; een cirkel heeft er één (zijn
middelpunt); een driehoek nooit.
\end{definition}

\section{Parallellogrammen}

\begin{definition}[Parallellogram]\label{def:g7:central:parallelogram}
Een \emph{parallellogram}\index{parallellogram} is een \omterm{def:g4:geometry:polygon}{vierhoek} waarvan de
diagonalen elkaar in hun gemeenschappelijke midden snijden. Dat snijpunt is dan een
\omterm{def:g7:central:center}{symmetriemiddelpunt} van de figuur: elk \omterm{def:g5:solids:def}{hoekpunt} is het beeld van het \omterm{def:g5:solids:def}{hoekpunt}
ertegenover onder de halve draai.
\end{definition}

\begin{theorem}[Eigenschappen van een parallellogram]\label{thm:g7:central:properties}
In een \omterm{def:g7:central:parallelogram}{parallellogram} $ABCD$ met middelpunt $O$:
\begin{enumerate}
  \item zijn de zijden die tegenover elkaar liggen \omterm{def:g6:lines:perp}{parallel}:
        $(AB) \parallel (CD)$ en $(BC) \parallel (AD)$;
  \item hebben de zijden die tegenover elkaar liggen dezelfde lengte: $AB = CD$ en
        $BC = AD$;
  \item zijn de hoeken die tegenover elkaar liggen gelijk: $\widehat A = \widehat C$
        en $\widehat B = \widehat D$.
\end{enumerate}
\end{theorem}

\begin{proof}
De halve draai om $O$ stuurt $A$ naar $C$ en $B$ naar $D$ (per definitie is $O$ het
midden van beide diagonalen). Ze stuurt dus het \omterm{def:g6:lines:objects}{lijnstuk} $[AB]$ naar het \omterm{def:g6:lines:objects}{lijnstuk}
$[CD]$; volgens \cref{prop:g7:central:preserved} is het beeld \omterm{def:g6:lines:perp}{parallel} aan $[AB]$
en even lang: punten 1 en 2. Ze stuurt ook de hoek in $A$ naar de hoek in $C$, en
hoeken blijven bewaard: punt 3.
\end{proof}

\begin{omfigure}
\begin{tikzpicture}[scale=1.0]
  \coordinate (A) at (0,0);
  \coordinate (B) at (3.4,0);
  \coordinate (C) at (4.6,1.9);
  \coordinate (D) at (1.2,1.9);
  \draw[thick] (A) -- (B) -- (C) -- (D) -- cycle;
  \draw[omThm] (A) -- (C);
  \draw[omThm] (B) -- (D);
  \fill (2.3,0.95) circle (1.6pt) node[right=3pt, font=\small] {$O$};
  \node[below left, font=\small] at (A) {$A$};
  \node[below right, font=\small] at (B) {$B$};
  \node[above right, font=\small] at (C) {$C$};
  \node[above left, font=\small] at (D) {$D$};
  \draw (1.62,-0.09) -- (1.78,0.09);
  \draw (2.82,1.81) -- (2.98,1.99);
  \draw (0.53,1.02) -- (0.67,0.88);
  \draw (3.93,1.02) -- (4.07,0.88);
\end{tikzpicture}

{\small Een \omterm{def:g7:central:parallelogram}{parallellogram} met zijn middelpunt $O$: de diagonalen (rood) snijden
elkaar in hun middens, en de zijden die tegenover elkaar liggen zijn \omterm{def:g6:lines:perp}{parallel} en
even lang (streepjes).}
\end{omfigure}

\begin{theorem}[Een parallellogram herkennen]\label{thm:g7:central:recognize}
Een \omterm{def:g4:geometry:polygon}{vierhoek} (zonder kruisende zijden) is al een \omterm{def:g7:central:parallelogram}{parallellogram} zodra \emph{één} van
de volgende dingen geldt:
\begin{enumerate}
  \item zijn diagonalen hebben hetzelfde midden (de definitie);
  \item de zijden die tegenover elkaar liggen zijn twee aan twee \omterm{def:g6:lines:perp}{parallel};
  \item twee zijden die tegenover elkaar liggen zijn \omterm{def:g6:lines:perp}{parallel} \emph{en} even lang.
\end{enumerate}
\end{theorem}
\admitted

\begin{example}\label{ex:g7:central:recognize}
Van $ABCD$ is gegeven dat $(AB) \parallel (CD)$ en $AB = CD = 5$ cm. Criterium 3
geldt: $ABCD$ is een \omterm{def:g7:central:parallelogram}{parallellogram} --- en dus weten we zonder verder te meten ook
dat $AD = BC$, dat $(AD) \parallel (BC)$, en dat zijn diagonalen elkaar middendoor
snijden.
\end{example}

\begin{remark}[Bijzondere parallellogrammen]
De rechthoek, de \omterm{def:g6:shapes:quadrilaterals}{ruit} en het vierkant van \cref{ch:g6:shapes} zijn precies de
parallellogrammen met een extra eigenschap: gelijke diagonalen (rechthoek),
loodrechte diagonalen (\omterm{def:g6:shapes:quadrilaterals}{ruit}), of beide (vierkant). Het hoekargument dat in
\cref{ex:g6:shapes:recognize} beloofd werd, werkt nu: in een \omterm{def:g7:central:parallelogram}{parallellogram} zijn
opeenvolgende hoeken supplementair (verwisselende en overeenkomstige hoeken bij de
parallellen, \cref{thm:g7:triangles:equalities}), dus dwingt één \omterm{def:g3:shapes:rightangle}{rechte hoek} alle
vier.
\end{remark}

\section{Oefeningen}

\begin{exercise}[$\star$]\label{exo:g7:central:1}
Zet op ruitjespapier $O$, en daarna $M$ vier vakjes rechts van $O$ en één vakje
hoger. Construeer het \omterm{def:g7:central:def}{spiegelbeeld} $M'$ van $M$ om $O$, en het \omterm{def:g7:central:def}{spiegelbeeld} van $M$
in de \emph{verticale rasterlijn} door $O$. Vergelijk de twee beelden.
\end{exercise}

\begin{exercise}[$\star$]\label{exo:g7:central:2}
Teken een driehoek $ABC$ en een punt $O$ erbuiten. Construeer zijn beeld $A'B'C'$
onder de halve draai om $O$. Wat kun je zeggen over de lengtes $A'B'$ en $AB$? En
over de \omterm{def:g6:lines:objects}{lijnen} $(A'B')$ en $(AB)$?
\end{exercise}

\begin{exercise}[$\star$]\label{exo:g7:central:3}
Welke cijfers, geschreven zoals op een rekenmachinedisplay ($0$ tot $9$), hebben een
\omterm{def:g7:central:center}{symmetriemiddelpunt}? En welke hoofdletters van H, A, N, S, T, Z, O?
\end{exercise}

\begin{exercise}[$\star$]\label{exo:g7:central:4}
$KLMN$ is een \omterm{def:g7:central:parallelogram}{parallellogram} met middelpunt $O$, waarbij $KL = 7$ cm, $LM = 4$ cm en
$\widehat K = 115^\circ$. Geef, met redenen: $MN$, $NK$, $\widehat M$, en het midden
van $[LN]$.
\end{exercise}

\begin{exercise}[$\star$]\label{exo:g7:central:5}
Construeer een \omterm{def:g7:central:parallelogram}{parallellogram} $ABCD$ uit zijn diagonalen: $AC = 8$ cm en
$BD = 5$ cm, die elkaar in hun middens snijden onder een hoek van $60^\circ$. (Teken
eerst de diagonalen!)
\end{exercise}

\begin{exercise}[$\star$]\label{exo:g7:central:6}
In een \omterm{def:g7:central:parallelogram}{parallellogram} meet één hoek $115^\circ$. Geef met de opmerking over
supplementaire opeenvolgende hoeken de drie andere hoeken.
\end{exercise}

\begin{exercise}[$\star\star$]\label{exo:g7:central:7}
Zet $A(1, 1)$, $B(4, 2)$ en het punt $O(2.5, 2.5)$ in een assenstelsel. Reken de
coördinaten uit van de beelden $A'$ en $B'$ onder de halve draai om $O$ (gebruik het
tellen op het raster van \cref{ex:g7:central:grid}).
\end{exercise}

\begin{exercise}[$\star\star$]\label{exo:g7:central:8}
$ABCD$ is een \omterm{def:g4:geometry:polygon}{vierhoek} waarin $AB = CD$. Is dat genoeg om er een \omterm{def:g7:central:parallelogram}{parallellogram} van
te maken? Zo niet, schets dan een tegenvoorbeeld (een \omterm{def:g6:shapes:triangles}{gelijkbenig} trapezium helpt) en
zeg wat er aan de veronderstelling toegevoegd moet worden.
\end{exercise}

\begin{exercise}[$\star\star$]\label{exo:g7:central:9}
Zij $ABC$ een driehoek en $O$ het midden van $[BC]$. Construeer het \omterm{def:g7:central:def}{spiegelbeeld}
$A'$ van $A$ om $O$. Laat zien dat $ABA'C$ een \omterm{def:g7:central:parallelogram}{parallellogram} is. (Welk criterium van
\cref{thm:g7:central:recognize} krijg je hier gratis?)
\end{exercise}

\begin{exercise}[$\star\star$]\label{exo:g7:central:10}
Waar of niet waar, met een reden: ``een \omterm{def:g4:geometry:polygon}{vierhoek} waarvan de hoeken die tegenover
elkaar liggen twee aan twee gelijk zijn, is een \omterm{def:g7:central:parallelogram}{parallellogram}''; ``een \omterm{def:g7:central:parallelogram}{parallellogram}
met loodrechte diagonalen is een \omterm{def:g6:shapes:quadrilaterals}{ruit}''; ``een \omterm{def:g7:central:parallelogram}{parallellogram} met één \omterm{def:g3:shapes:rightangle}{rechte hoek} is
een rechthoek''.
\end{exercise}

\begin{exercise}[$\star\star\star$]\label{exo:g7:central:11}
Zij $ABCD$ een willekeurige \omterm{def:g4:geometry:polygon}{vierhoek} en $I$, $J$, $K$, $L$ de middens van zijn
zijden $[AB]$, $[BC]$, $[CD]$ en $[DA]$. Teken verschillende voorbeelden. Wat lijkt
$IJKL$ altijd te zijn? (Het bewijs, met de middenparallelstelling, komt in
\cref{ch:g8:midpoints} --- en nog eens met vectoren, in het bovenbouwvolume.)
\end{exercise}

\section{Opgave: de cake, de zwaartelijn en het onmogelijke tweede middelpunt}

\begin{problem}[{Weekendopgave --- halve draaien aan het werk: eerlijke
cakesneden, een grens voor zwaartelijnen, en waarom geen figuur twee
symmetriemiddelpunten kan hebben}]
\label{pb:g7:central:1}
De halve draai lijkt de eenvoudigste van alle transformaties --- en toch snijdt hij
cakes in bewijsbaar eerlijke \omterm{def:g3:division:half}{helften}, meet hij de zwaartelijnen van een driehoek, en
bergt hij een pareltje pure wiskunde: een begrensde figuur kan één
\omterm{def:g7:central:center}{symmetriemiddelpunt} hebben, maar \emph{nooit twee}. Het bewijs gebruikt een
ontdekking die je in deel I doet: twee halve draaien, de een na de ander, komen samen
neer op een verschuiving.

\medskip
\noindent\textbf{Deel I --- Gymnastiek met halve draaien.}
\begin{enumerate}
  \item Neem op ruitjespapier de oorsprong $O(0, 0)$ als middelpunt. Geef met het
        tellen op het raster van \cref{ex:g7:central:grid} de beelden van
        $M(3, 1)$, $N(-2, 4)$ en $P(0, -5)$ onder de halve draai om $O$. Wat is de
        algemene regel voor $(x, y)$?
  \item Nu is het middelpunt $C(2, 1)$. Reken het beeld van $M(5, 3)$ uit, en
        controleer het algemene recept: elke coördinaat van het beeld is
        \emph{twee keer die van het middelpunt min die van het punt}.
  \item Speelkaarten zijn zo ontworpen dat ze op hun kop hetzelfde lezen --- een
        \omterm{def:g7:central:def}{puntspiegeling}, zodat geen van beide spelers de kaart ``verkeerd'' ziet. Leg
        uit waarom op een kaart met een \emph{\omterm{def:g3:numbers:evenodd}{oneven}} aantal symbolen één symbool
        precies in het middelpunt van de kaart moet liggen.
  \item Twee halve draaien op een rij: neem de middelpunten $O_1(2, 0)$ en
        $O_2(5, 0)$. Stuur het punt $A(1, 1)$ door de halve draai om $O_1$, en stuur
        het beeld daarna door de halve draai om $O_2$. Vergelijk begin en einde.
        Herhaal met $B(0, 3)$. Op welke enkele, eenvoudige transformatie kwamen de
        twee halve draaien samen neer?
  \item Meet de verschuiving van vraag 4 af tegen de afstand $O_1 O_2$ en formuleer
        de ontdekking. (Vergelijk met de twee parallelle spiegels van
        \cref{pb:g6:symmetry:1}: hetzelfde verschijnsel, nieuwe spelers.)
\end{enumerate}

\medskip
\noindent\textbf{Deel II --- Cakes en zwaartelijnen.}
\begin{enumerate}[resume]
  \item Zij $O$ het middelpunt van een \omterm{def:g7:central:parallelogram}{parallellogram}. Leg uit waarom elke \omterm{def:g6:lines:objects}{lijn} door
        $O$ de rand in twee punten snijdt die \omterm{def:g7:central:def}{spiegelbeelden} om $O$ zijn, en waarom de
        \omterm{def:g6:lines:objects}{lijn} het \omterm{def:g7:central:parallelogram}{parallellogram} in twee stukken verdeelt die precies elkaars beeld
        onder de halve draai zijn --- en dus dezelfde \omterm{def:g6:measure:area}{oppervlakte} hebben.
  \item De stelling van de bakker: een rechthoekige cake wordt volmaakt eerlijk
        verdeeld door \emph{elke} rechte snede door zijn middelpunt. Sterker nog: een
        rechthoekige cake heeft een \omterm{def:g6:shapes:triangles}{rechthoekig} gat (elke maat, elke plaats, zelfs
        schuin). Beschrijf de ene rechte snede die twee delen met even veel cake geeft
        --- en verantwoord haar met vraag 6.
  \item Drie \omterm{def:g5:solids:def}{hoekpunten} van een \omterm{def:g7:central:parallelogram}{parallellogram} $ABCD$ zijn $A(0, 0)$, $B(6, 1)$ en
        $C(8, 5)$. Reken de coördinaten van $D$ en van het middelpunt $O$ uit.
  \item \Cref{exo:g7:central:9} bouwde uit een driehoek $ABC$ en het midden $O$ van
        $[BC]$ het \omterm{def:g7:central:parallelogram}{parallellogram} $ABA'C$, waarbij $A'$ het \omterm{def:g7:central:def}{spiegelbeeld} van $A$ om
        $O$ is. Gebruik dat, samen met de driehoeksongelijkheid
        (\cref{thm:g7:triangles:inequality}) in de driehoek $ABA'$, om de
        \emph{grens voor de zwaartelijn} te bewijzen: de zwaartelijn $[AO]$ voldoet
        aan
        \[
        AO < \frac{AB + AC}{2} .
        \]
        (Merk op dat $AA' = 2\,AO$ en $BA' = AC$.)
  \item Tussen welke twee waarden moet in een driehoek met $AB = 5$ cm en
        $AC = 7$ cm de lengte van de zwaartelijn uit $A$ liggen? (Gebruik beide
        richtingen van de driehoeksongelijkheid in $ABA'$.)
\end{enumerate}

\medskip
\noindent\textbf{Deel III --- Hoeveel middelpunten kan een figuur hebben?}
\begin{enumerate}[resume]
  \item Zeg van elke figuur of ze een \omterm{def:g7:central:center}{symmetriemiddelpunt} heeft, en waar: een
        \omterm{def:g6:lines:objects}{lijnstuk}; een volledige \omterm{def:g6:lines:objects}{lijn}; een \omterm{def:g6:shapes:triangles}{gelijkzijdige driehoek}; een vierkant; een
        cirkel; de letters S en Z.
  \item Bewijs je bewering voor de \omterm{def:g6:shapes:triangles}{gelijkzijdige driehoek}: een halve draai die de
        driehoek op zichzelf afbeeldt, zou de drie \omterm{def:g5:solids:def}{hoekpunten} in paren zetten ---
        maar drie is \omterm{def:g3:numbers:evenodd}{oneven}, dus zou een \omterm{def:g5:solids:def}{hoekpunt} naar zichzelf gestuurd worden. Wat
        zou dat \omterm{def:g5:solids:def}{hoekpunt} dan moeten zijn, en waarom is dat onmogelijk?
  \item Veralgemeen vraag 12: leg uit waarom \emph{geen} \omterm{def:g4:geometry:polygon}{veelhoek} met een \omterm{def:g3:numbers:evenodd}{oneven}
        aantal \omterm{def:g5:solids:def}{hoekpunten} een \omterm{def:g7:central:center}{symmetriemiddelpunt} heeft.
  \item De parel: stel dat een figuur \emph{twee} verschillende
        \omterm{def:g7:central:center}{symmetriemiddelpunten} $O_1$ en $O_2$ had. Voer de twee halve draaien na
        elkaar uit en gebruik de ontdekking van vraag 5: welke beweging zou de figuur
        dan precies op zichzelf moeten afbeelden? Leg uit waarom een figuur die op een
        blad papier past dat niet kan overleven, en besluit.
  \item De conclusie van vraag 14 laat een achterdeurtje open: \emph{onbegrensde}
        figuren. Ga op een oneindige fries --- bijvoorbeeld een eindeloze strook
        voetstappen, links, rechts, links, rechts\,\dots\ --- na dat ze (minstens)
        twee verschillende \omterm{def:g7:central:center}{symmetriemiddelpunten} heeft, en wijs de verschuiving aan
        die vraag 14 voorspelde.
\end{enumerate}
\end{problem}
